% Methods subsections --- Data + Conformal Statistics (Article 1).
% Team A, Week 5 Day 31. Source of record: week5/reports/w5_09_methods_data_stats.md
% Drop into the Methods section with % Methods subsections --- Data + Conformal Statistics (Article 1).
% Team A, Week 5 Day 31. Source of record: week5/reports/w5_09_methods_data_stats.md
% Drop into the Methods section with % Methods subsections --- Data + Conformal Statistics (Article 1).
% Team A, Week 5 Day 31. Source of record: week5/reports/w5_09_methods_data_stats.md
% Drop into the Methods section with % Methods subsections --- Data + Conformal Statistics (Article 1).
% Team A, Week 5 Day 31. Source of record: week5/reports/w5_09_methods_data_stats.md
% Drop into the Methods section with \input{w5_09_methods}.
% Requires: \usepackage{booktabs}, \usepackage{amssymb}, \usepackage{amsmath}.
% Cells marked \ddag are provisional (Day-14 dummy fidelity interface) and reprice on real prototypes.

\subsection{Data}
\label{sec:methods-data}

\paragraph{Datasets and unified representation.}
We evaluate on three public IoT/network intrusion-detection corpora: CIC-IoT2023, BoT-IoT and UNSW-NB15.
Two further corpora (TON\_IoT, Edge-IIoTset) were curated under the same pipeline but excluded, because
their held-out zero-day split does not survive the shared schema: projecting onto the common feature set
maps many distinct attack flows onto \emph{identical} vectors, and the cross-split de-duplication that
follows resolves collisions with training-split priority, so the surviving zero-day copies are dropped. The
held-out families end at one and two rows respectively, which cannot support a rejection-rate estimate; the
three retained datasets keep their zero-day splits intact. Since the three corpora publish incompatible
column vocabularies,
each is projected onto a shared \textbf{17-feature unified flow schema} (flow duration; total, source and
destination packet and byte counts; overall, source and destination rates; four packet-size statistics;
inter-arrival time; and harmonized \texttt{protocol} and \texttt{conn\_state} categoricals), together with a
fine multiclass label, a binary label, and a ten-family attack-ontology label.

\paragraph{Curation and leakage controls.}
Curation runs through a quantum-aware curation pipeline whose defining property is that it is
\emph{split-first}: every data-dependent parameter is fit on the training split alone and only applied to
the others. This departs deliberately from the conventional ordering, in which scaling or balancing precedes
splitting and thereby leaks evaluation data into fitted parameters. Concretely, the correlation prune
(dropping one of each $|r| \geq 0.95$ pair), the median/IQR robust scalers, the angle-encoding ranges
(mapping features to $[0, \pi]$ for angle embedding; the published benchmark uses the qubit budget of $8$
from the nested $4/8/12/16$ rankings), and the feature ranking are all fit on the training split and applied
unchanged to calibration, test and zero-day. Class balancing is applied to
training only; evaluation splits are never resampled.

Two further controls underpin the guarantee developed in Section~\ref{sec:methods-conformal}. First,
de-duplication is performed on the feature columns \emph{before} splitting, so one feature vector can land
in exactly one split; cross-split disjointness is then asserted on feature-content hashes rather than row
indices, which is what detects a repeated flow shared across splits. This is the relevant control for
CIC-IoT2023, whose leakage vector is repeated and near-duplicate flows rather than identifier columns.
Measured cross-split flow overlap is zero for every pair on every dataset, and within-split duplication is
zero. Second, the conformal threshold is calibrated on held-out \emph{known-class} data and never on the
unknown class: calibrating on the held-out family would leak precisely the signal the detector is meant to
discover.

\paragraph{Partition protocol.}
Each dataset is partitioned four ways --- Train / Calibration / Test / Zero-Day --- under a fixed seed
($42$), so the partitions are deterministic and reproduce exactly on re-run. The zero-day split is a
\emph{held-out attack family} rather than a random subsample: one or more entire attack classes are removed
from training, calibration and test and appear only at evaluation time. The remaining classes are stratified
$70/10/10/10$ into train/validation/calibration/test, with validation folded into train to give the
four-way conformal layout. Calibration contains known classes only, by construction. Three invariants are
asserted in the partitioner rather than merely documented: the zero-day family has zero rows outside the
zero-day split; calibration is a subset of the known classes; and the assignment is deterministic under the
seed.

\begin{table}[t]
\centering
\caption{Partition protocol. Row counts per split, number of known classes, and the held-out zero-day
family. Seed 42; deterministic on re-run.}
\label{tab:partitions}
\begin{tabular}{lrrrrrl}
\toprule
Dataset & Train & Calib. & Test & Zero-day & \#known & Held-out family \\
\midrule
CIC-IoT2023 & 151,049 & 18,883 & 18,883 & 10,984 & 31 & Mirai ($\times 3$) \\
BoT-IoT     & 148,729 & 18,591 & 18,591 &    683 &  4 & Theft \\
UNSW-NB15   &  81,215 & 10,112 & 10,086 &  1,216 &  8 & Shellcode, Worms \\
\bottomrule
\end{tabular}
\end{table}

\paragraph{Verifying exchangeability.}
Split-conformal validity requires \emph{exchangeability} between the calibration set and each test point ---
strictly weaker than i.i.d., and a condition that cannot be proven, only falsified. We therefore treat it as
an empirical claim and attempt to falsify it with a classifier two-sample test: a random forest is trained
to separate calibration rows from test rows and scored by five-fold cross-validated AUROC, where chance
performance means the two sets are indistinguishable on the features. We obtain AUROC $0.4999$, $0.5038$ and
$0.5035$ on CIC-IoT2023, BoT-IoT and UNSW-NB15 respectively, with cal-vs-test total-variation distances of
$0.0000$, $0.0000$ and $0.0019$; exchangeability is not rejected on any dataset. This claim is scoped to the
\emph{known-class} calibration--test split. The held-out zero-day split is deliberately non-exchangeable
with calibration --- that non-exchangeability is the detection signal --- so the guarantee below covers the
known-class portion of the test stream, and zero-day detection is reported separately.

\subsection{Conformal Statistics}
\label{sec:methods-conformal}

\paragraph{Nonconformity score and calibration.}
For a sample $x$ encoded as a density matrix $\rho_x$, the model supplies the Uhlmann fidelity
$F(\rho_x, \rho_c)$ to every known-class prototype $\rho_c$ (non-squared $F$). The nonconformity score is
\begin{equation}
s(x) \;=\; 1 - \max_{c} F(\rho_x, \rho_c),
\end{equation}
so a sample with low fidelity to every known prototype scores high and is more novel. Prototypes are built
from the training split only; were calibration rows to contribute to $\rho_c$, calibration scores would be
in-sample and the threshold optimistic. The threshold is the split-conformal order statistic of $s$ over the
known-class calibration set of size $n$,
\begin{equation}
k = \left\lceil (1-\alpha)(n+1) \right\rceil, \qquad q = s_{(k)},
\end{equation}
the $k$-th smallest calibration score. A test point is flagged zero-day iff $s > q$. Under exchangeability
of the known-class calibration and test points this yields the finite-sample marginal guarantee
\begin{equation}
1 - \alpha \;\leq\; \mathbb{P}(s_{\mathrm{test}} \leq q) \;\leq\; 1 - \alpha + \frac{1}{n+1},
\label{eq:coverage}
\end{equation}
i.e.\ the false-zero-day rate on known traffic is at most $\alpha$, with granularity $1/(n+1)$. Validity is
\emph{score-agnostic}: it does not require the fidelity score to be a good novelty signal. Detection
\emph{power} does, and is therefore reported as an empirical result rather than inherited from
Eq.~\eqref{eq:coverage} --- a coverage curve is not a power curve. We use one primary $\alpha = 0.05$ for
the quantum and classical arms alike, giving $k = 17{,}940$, $17{,}663$ and $9{,}608$ and thresholds
$q = 0.3006^{\ddag}$, $0.2629^{\ddag}$ and $0.3834^{\ddag}$ on the three datasets.

\paragraph{Marginal rather than class-conditional calibration.}
Class-conditional (Mondrian) thresholds are degenerate when a class holds fewer than
$\lceil 1/\alpha \rceil - 1$ calibration points, a floor of $19$ per class at $\alpha = 0.05$. CIC-IoT2023
has four known classes with at most nine calibration rows, so we calibrate marginally on all three datasets
for comparability and report class-conditional behaviour as a diagnostic. The scope of the marginal
guarantee is stated precisely, because it is easy to overclaim: marginal conformal controls the false-alarm
rate over the known-class \emph{mixture} and carries \emph{no per-class guarantee}. Rare classes contribute
few pooled calibration points and so do not skew $q$ --- the mixture guarantee is unaffected by class
imbalance --- but an individual class may still be over- or under-covered. Each known class's test flag
count is therefore judged against its own exact band; on the present interface, $0$ of $31$, $0$ of $4$ and
$0$ of $8$ classes fall outside their own band. Should a class fall outside on real prototypes, the
appropriate remedy is clustered conformal --- grouping rare classes until each cluster clears the
calibration floor --- rather than fully class-conditional conformal, which several classes cannot support.

\paragraph{Coverage verification.}
A bare $\mathrm{FZR} \leq \alpha$ assertion is the wrong acceptance test. Split conformal guarantees the
\emph{expectation}, marginally over the calibration and test draws; on a finite test set the empirical rate
fluctuates around $(n+1-k)/(n+1)$, so a naive threshold check rejects sound systems by luck roughly half the
time. We instead judge each observation against the exact finite-sample law: conditional on the calibration
draw, coverage is Beta-distributed, so the number of false flags $E$ on $m$ known test rows follows a
Beta-Binomial,
\begin{equation}
\mathrm{coverage} \mid \mathrm{cal} \sim \mathrm{Beta}(k,\, n{+}1{-}k)
\quad \Longrightarrow \quad
E \sim \mathrm{BetaBin}(m,\, n{+}1{-}k,\, k).
\end{equation}
For each (dataset, $\alpha$) we report the observed rate, its exact central $99\%$ band, and a one-sided
tail $p$-value $\mathbb{P}(E \geq e_{\mathrm{obs}})$. The band, not the point comparison against $\alpha$,
is the acceptance gate: a failure means the guarantee is genuinely violated rather than unlucky. Below-band
excursions are reported as conservativeness --- ties and score discreteness make the band conservative ---
not as failures. This is also the correct band for a calibration-draw plot; Clopper--Pearson describes a
fixed-threshold proportion and is the wrong object here. BoT-IoT illustrates the distinction: its achieved
rate $0.0530^{\ddag}$ exceeds $\alpha$ outright, yet sits comfortably inside the exact band
$[0.0443, 0.0559]$ at $p = 0.090$, which is ordinary sampling noise on $18{,}591$ test rows. Coverage is
verified across the full sweep $\alpha \in [0.01, 0.20]$ ($20$ values $\times$ $3$ datasets $=60$
verifications, all passing), and the achieved rate tracks the diagonal $y = \alpha$ inside the band across
the whole range.

\paragraph{Stability across seeds.}
Because the guarantee is marginal over the calibration/test draw, it must hold across draws rather than on
one split. We pool the known-class calibration and test scores, re-draw them at the original sizes under
seeds $42$--$46$, re-calibrate each re-split and judge it against its own exact band, holding the zero-day
pool fixed and scoring it at each re-split's $q$. All three datasets hold the exact band on all five
re-splits, and every $95\%$ confidence interval covers the target $\alpha = 0.05$ (mean achieved rates
$0.0517^{\ddag}$, $0.0501^{\ddag}$, $0.0508^{\ddag}$). The accompanying recall intervals quantify how little
the power number moves with the calibration draw.

\paragraph{Auditing the assumption at deployment.}
Coverage arithmetic is meaningful only if the assumption underneath it holds, so the deployed pipeline is
audited with six inferential tests per dataset --- two-sample Kolmogorov--Smirnov and Mann--Whitney on
calibration versus test scores, a class-mix $\chi^2$, a flag-rate-by-class $\chi^2$, and two index-drift
Spearman tests --- with Holm step-down over the declared family of $18$ cells. No violations are detected.
Two choices deserve statement. First, the textbook conformal $p$-value uniformity check is reported
descriptively and \emph{excluded} from the inferential family: every $p_i$ is computed against the same
calibration set, so the $p$-values are exchangeable rather than independent and the one-sample uniformity
null is far too tight (measured on this interface it rejects on approximately $37\%$ of genuinely
exchangeable re-splits); it is also redundant with the two-sample statistic up to $1/(n+1)$. Second, the
audit's power is demonstrated by injected-violation drills rather than assumed: trimming the calibration
tail, shifting test scores, and skewing the test class mix are each injected and each detected. The
class-mix injection is instructive --- the audit flags it decisively while the \emph{marginal} false-alarm
rate stays inside the band, which is exactly the failure a coverage-only check cannot see. A caveat travels
with the remedy: re-splitting repairs \emph{split-induced} non-exchangeability, and also makes a
\emph{deployment shift} look repaired, because calibration and test become exchangeable draws from the
shifted mixture --- but that mixture is not the deployment distribution. A genuine deployment shift is
answered by recalibrating on fresh data, never by re-splitting.

\paragraph{Significance protocol.}
All comparative claims pass through one pre-declared protocol, chosen so the test matches the object being
claimed. We report mean $\pm$ sample standard deviation (ddof $=1$) with $t$-based $95\%$ confidence
intervals. Seed-level comparisons use a two-sided paired $t$-test with pairs matched by seed over seeds
$42$--$46$; row-level comparisons use an exact two-sided binomial McNemar test on discordant pairs matched
by test row. Effect sizes accompany every $p$-value: paired $d_z = \mathrm{mean}(\Delta)/\mathrm{sd}(\Delta)$
for seed-level tests and Cohen's $h$ for proportions. Multiplicity is controlled by Holm--Bonferroni
step-down over a family declared \emph{before} testing, and a zero-variance guard forces $t = 0$, $p = 1$,
$d = 0$ so a self-comparison cannot fake significance. Seed-paired and row-paired tests answer different
questions and are reported separately rather than pooled: threshold-free ranking metrics speak to ordering,
whereas McNemar tests the flags a system actually deploys. A dry run on classical baselines showed the
correction earning its keep, with a raw $p = 0.011$ in a nine-test family correctly losing significance
under Holm at $p = 0.056$.

Families are declared per comparison surface: closed-set correctness against the classical multiclass
detector ($m = 3$, one comparison per dataset) and known-split false-alarm rate against every novelty head
($m = 8$; one dataset publishes no one-class SVM model, a real absence rather than a missing value),
corrected separately. We note explicitly that the false-alarm family compares systems at their
\emph{shipped operating points} --- the conformal detector at its guaranteed $\alpha = 0.05$, each
heuristic head at its own frozen threshold, which targeted a $2\alpha$ known-FPR budget --- so it supports a
deployed-behaviour statement rather than a matched-budget superiority claim; the like-for-like comparison at
one shared $\alpha$ is reported with the zero-day results. Finally, the seed-level power floor is stated
rather than left implicit: with $n = 5$ seeds the smallest detectable paired effect is $d_z \approx 1.24$,
and $\approx 1.68$ at $80\%$ power, which is why an effect size accompanies every reported $p$-value.
.
% Requires: \usepackage{booktabs}, \usepackage{amssymb}, \usepackage{amsmath}.
% Cells marked \ddag are provisional (Day-14 dummy fidelity interface) and reprice on real prototypes.

\subsection{Data}
\label{sec:methods-data}

\paragraph{Datasets and unified representation.}
We evaluate on three public IoT/network intrusion-detection corpora: CIC-IoT2023, BoT-IoT and UNSW-NB15.
Two further corpora (TON\_IoT, Edge-IIoTset) were curated under the same pipeline but excluded, because
their held-out zero-day split does not survive the shared schema: projecting onto the common feature set
maps many distinct attack flows onto \emph{identical} vectors, and the cross-split de-duplication that
follows resolves collisions with training-split priority, so the surviving zero-day copies are dropped. The
held-out families end at one and two rows respectively, which cannot support a rejection-rate estimate; the
three retained datasets keep their zero-day splits intact. Since the three corpora publish incompatible
column vocabularies,
each is projected onto a shared \textbf{17-feature unified flow schema} (flow duration; total, source and
destination packet and byte counts; overall, source and destination rates; four packet-size statistics;
inter-arrival time; and harmonized \texttt{protocol} and \texttt{conn\_state} categoricals), together with a
fine multiclass label, a binary label, and a ten-family attack-ontology label.

\paragraph{Curation and leakage controls.}
Curation runs through a quantum-aware curation pipeline whose defining property is that it is
\emph{split-first}: every data-dependent parameter is fit on the training split alone and only applied to
the others. This departs deliberately from the conventional ordering, in which scaling or balancing precedes
splitting and thereby leaks evaluation data into fitted parameters. Concretely, the correlation prune
(dropping one of each $|r| \geq 0.95$ pair), the median/IQR robust scalers, the angle-encoding ranges
(mapping features to $[0, \pi]$ for angle embedding; the published benchmark uses the qubit budget of $8$
from the nested $4/8/12/16$ rankings), and the feature ranking are all fit on the training split and applied
unchanged to calibration, test and zero-day. Class balancing is applied to
training only; evaluation splits are never resampled.

Two further controls underpin the guarantee developed in Section~\ref{sec:methods-conformal}. First,
de-duplication is performed on the feature columns \emph{before} splitting, so one feature vector can land
in exactly one split; cross-split disjointness is then asserted on feature-content hashes rather than row
indices, which is what detects a repeated flow shared across splits. This is the relevant control for
CIC-IoT2023, whose leakage vector is repeated and near-duplicate flows rather than identifier columns.
Measured cross-split flow overlap is zero for every pair on every dataset, and within-split duplication is
zero. Second, the conformal threshold is calibrated on held-out \emph{known-class} data and never on the
unknown class: calibrating on the held-out family would leak precisely the signal the detector is meant to
discover.

\paragraph{Partition protocol.}
Each dataset is partitioned four ways --- Train / Calibration / Test / Zero-Day --- under a fixed seed
($42$), so the partitions are deterministic and reproduce exactly on re-run. The zero-day split is a
\emph{held-out attack family} rather than a random subsample: one or more entire attack classes are removed
from training, calibration and test and appear only at evaluation time. The remaining classes are stratified
$70/10/10/10$ into train/validation/calibration/test, with validation folded into train to give the
four-way conformal layout. Calibration contains known classes only, by construction. Three invariants are
asserted in the partitioner rather than merely documented: the zero-day family has zero rows outside the
zero-day split; calibration is a subset of the known classes; and the assignment is deterministic under the
seed.

\begin{table}[t]
\centering
\caption{Partition protocol. Row counts per split, number of known classes, and the held-out zero-day
family. Seed 42; deterministic on re-run.}
\label{tab:partitions}
\begin{tabular}{lrrrrrl}
\toprule
Dataset & Train & Calib. & Test & Zero-day & \#known & Held-out family \\
\midrule
CIC-IoT2023 & 151,049 & 18,883 & 18,883 & 10,984 & 31 & Mirai ($\times 3$) \\
BoT-IoT     & 148,729 & 18,591 & 18,591 &    683 &  4 & Theft \\
UNSW-NB15   &  81,215 & 10,112 & 10,086 &  1,216 &  8 & Shellcode, Worms \\
\bottomrule
\end{tabular}
\end{table}

\paragraph{Verifying exchangeability.}
Split-conformal validity requires \emph{exchangeability} between the calibration set and each test point ---
strictly weaker than i.i.d., and a condition that cannot be proven, only falsified. We therefore treat it as
an empirical claim and attempt to falsify it with a classifier two-sample test: a random forest is trained
to separate calibration rows from test rows and scored by five-fold cross-validated AUROC, where chance
performance means the two sets are indistinguishable on the features. We obtain AUROC $0.4999$, $0.5038$ and
$0.5035$ on CIC-IoT2023, BoT-IoT and UNSW-NB15 respectively, with cal-vs-test total-variation distances of
$0.0000$, $0.0000$ and $0.0019$; exchangeability is not rejected on any dataset. This claim is scoped to the
\emph{known-class} calibration--test split. The held-out zero-day split is deliberately non-exchangeable
with calibration --- that non-exchangeability is the detection signal --- so the guarantee below covers the
known-class portion of the test stream, and zero-day detection is reported separately.

\subsection{Conformal Statistics}
\label{sec:methods-conformal}

\paragraph{Nonconformity score and calibration.}
For a sample $x$ encoded as a density matrix $\rho_x$, the model supplies the Uhlmann fidelity
$F(\rho_x, \rho_c)$ to every known-class prototype $\rho_c$ (non-squared $F$). The nonconformity score is
\begin{equation}
s(x) \;=\; 1 - \max_{c} F(\rho_x, \rho_c),
\end{equation}
so a sample with low fidelity to every known prototype scores high and is more novel. Prototypes are built
from the training split only; were calibration rows to contribute to $\rho_c$, calibration scores would be
in-sample and the threshold optimistic. The threshold is the split-conformal order statistic of $s$ over the
known-class calibration set of size $n$,
\begin{equation}
k = \left\lceil (1-\alpha)(n+1) \right\rceil, \qquad q = s_{(k)},
\end{equation}
the $k$-th smallest calibration score. A test point is flagged zero-day iff $s > q$. Under exchangeability
of the known-class calibration and test points this yields the finite-sample marginal guarantee
\begin{equation}
1 - \alpha \;\leq\; \mathbb{P}(s_{\mathrm{test}} \leq q) \;\leq\; 1 - \alpha + \frac{1}{n+1},
\label{eq:coverage}
\end{equation}
i.e.\ the false-zero-day rate on known traffic is at most $\alpha$, with granularity $1/(n+1)$. Validity is
\emph{score-agnostic}: it does not require the fidelity score to be a good novelty signal. Detection
\emph{power} does, and is therefore reported as an empirical result rather than inherited from
Eq.~\eqref{eq:coverage} --- a coverage curve is not a power curve. We use one primary $\alpha = 0.05$ for
the quantum and classical arms alike, giving $k = 17{,}940$, $17{,}663$ and $9{,}608$ and thresholds
$q = 0.3006^{\ddag}$, $0.2629^{\ddag}$ and $0.3834^{\ddag}$ on the three datasets.

\paragraph{Marginal rather than class-conditional calibration.}
Class-conditional (Mondrian) thresholds are degenerate when a class holds fewer than
$\lceil 1/\alpha \rceil - 1$ calibration points, a floor of $19$ per class at $\alpha = 0.05$. CIC-IoT2023
has four known classes with at most nine calibration rows, so we calibrate marginally on all three datasets
for comparability and report class-conditional behaviour as a diagnostic. The scope of the marginal
guarantee is stated precisely, because it is easy to overclaim: marginal conformal controls the false-alarm
rate over the known-class \emph{mixture} and carries \emph{no per-class guarantee}. Rare classes contribute
few pooled calibration points and so do not skew $q$ --- the mixture guarantee is unaffected by class
imbalance --- but an individual class may still be over- or under-covered. Each known class's test flag
count is therefore judged against its own exact band; on the present interface, $0$ of $31$, $0$ of $4$ and
$0$ of $8$ classes fall outside their own band. Should a class fall outside on real prototypes, the
appropriate remedy is clustered conformal --- grouping rare classes until each cluster clears the
calibration floor --- rather than fully class-conditional conformal, which several classes cannot support.

\paragraph{Coverage verification.}
A bare $\mathrm{FZR} \leq \alpha$ assertion is the wrong acceptance test. Split conformal guarantees the
\emph{expectation}, marginally over the calibration and test draws; on a finite test set the empirical rate
fluctuates around $(n+1-k)/(n+1)$, so a naive threshold check rejects sound systems by luck roughly half the
time. We instead judge each observation against the exact finite-sample law: conditional on the calibration
draw, coverage is Beta-distributed, so the number of false flags $E$ on $m$ known test rows follows a
Beta-Binomial,
\begin{equation}
\mathrm{coverage} \mid \mathrm{cal} \sim \mathrm{Beta}(k,\, n{+}1{-}k)
\quad \Longrightarrow \quad
E \sim \mathrm{BetaBin}(m,\, n{+}1{-}k,\, k).
\end{equation}
For each (dataset, $\alpha$) we report the observed rate, its exact central $99\%$ band, and a one-sided
tail $p$-value $\mathbb{P}(E \geq e_{\mathrm{obs}})$. The band, not the point comparison against $\alpha$,
is the acceptance gate: a failure means the guarantee is genuinely violated rather than unlucky. Below-band
excursions are reported as conservativeness --- ties and score discreteness make the band conservative ---
not as failures. This is also the correct band for a calibration-draw plot; Clopper--Pearson describes a
fixed-threshold proportion and is the wrong object here. BoT-IoT illustrates the distinction: its achieved
rate $0.0530^{\ddag}$ exceeds $\alpha$ outright, yet sits comfortably inside the exact band
$[0.0443, 0.0559]$ at $p = 0.090$, which is ordinary sampling noise on $18{,}591$ test rows. Coverage is
verified across the full sweep $\alpha \in [0.01, 0.20]$ ($20$ values $\times$ $3$ datasets $=60$
verifications, all passing), and the achieved rate tracks the diagonal $y = \alpha$ inside the band across
the whole range.

\paragraph{Stability across seeds.}
Because the guarantee is marginal over the calibration/test draw, it must hold across draws rather than on
one split. We pool the known-class calibration and test scores, re-draw them at the original sizes under
seeds $42$--$46$, re-calibrate each re-split and judge it against its own exact band, holding the zero-day
pool fixed and scoring it at each re-split's $q$. All three datasets hold the exact band on all five
re-splits, and every $95\%$ confidence interval covers the target $\alpha = 0.05$ (mean achieved rates
$0.0517^{\ddag}$, $0.0501^{\ddag}$, $0.0508^{\ddag}$). The accompanying recall intervals quantify how little
the power number moves with the calibration draw.

\paragraph{Auditing the assumption at deployment.}
Coverage arithmetic is meaningful only if the assumption underneath it holds, so the deployed pipeline is
audited with six inferential tests per dataset --- two-sample Kolmogorov--Smirnov and Mann--Whitney on
calibration versus test scores, a class-mix $\chi^2$, a flag-rate-by-class $\chi^2$, and two index-drift
Spearman tests --- with Holm step-down over the declared family of $18$ cells. No violations are detected.
Two choices deserve statement. First, the textbook conformal $p$-value uniformity check is reported
descriptively and \emph{excluded} from the inferential family: every $p_i$ is computed against the same
calibration set, so the $p$-values are exchangeable rather than independent and the one-sample uniformity
null is far too tight (measured on this interface it rejects on approximately $37\%$ of genuinely
exchangeable re-splits); it is also redundant with the two-sample statistic up to $1/(n+1)$. Second, the
audit's power is demonstrated by injected-violation drills rather than assumed: trimming the calibration
tail, shifting test scores, and skewing the test class mix are each injected and each detected. The
class-mix injection is instructive --- the audit flags it decisively while the \emph{marginal} false-alarm
rate stays inside the band, which is exactly the failure a coverage-only check cannot see. A caveat travels
with the remedy: re-splitting repairs \emph{split-induced} non-exchangeability, and also makes a
\emph{deployment shift} look repaired, because calibration and test become exchangeable draws from the
shifted mixture --- but that mixture is not the deployment distribution. A genuine deployment shift is
answered by recalibrating on fresh data, never by re-splitting.

\paragraph{Significance protocol.}
All comparative claims pass through one pre-declared protocol, chosen so the test matches the object being
claimed. We report mean $\pm$ sample standard deviation (ddof $=1$) with $t$-based $95\%$ confidence
intervals. Seed-level comparisons use a two-sided paired $t$-test with pairs matched by seed over seeds
$42$--$46$; row-level comparisons use an exact two-sided binomial McNemar test on discordant pairs matched
by test row. Effect sizes accompany every $p$-value: paired $d_z = \mathrm{mean}(\Delta)/\mathrm{sd}(\Delta)$
for seed-level tests and Cohen's $h$ for proportions. Multiplicity is controlled by Holm--Bonferroni
step-down over a family declared \emph{before} testing, and a zero-variance guard forces $t = 0$, $p = 1$,
$d = 0$ so a self-comparison cannot fake significance. Seed-paired and row-paired tests answer different
questions and are reported separately rather than pooled: threshold-free ranking metrics speak to ordering,
whereas McNemar tests the flags a system actually deploys. A dry run on classical baselines showed the
correction earning its keep, with a raw $p = 0.011$ in a nine-test family correctly losing significance
under Holm at $p = 0.056$.

Families are declared per comparison surface: closed-set correctness against the classical multiclass
detector ($m = 3$, one comparison per dataset) and known-split false-alarm rate against every novelty head
($m = 8$; one dataset publishes no one-class SVM model, a real absence rather than a missing value),
corrected separately. We note explicitly that the false-alarm family compares systems at their
\emph{shipped operating points} --- the conformal detector at its guaranteed $\alpha = 0.05$, each
heuristic head at its own frozen threshold, which targeted a $2\alpha$ known-FPR budget --- so it supports a
deployed-behaviour statement rather than a matched-budget superiority claim; the like-for-like comparison at
one shared $\alpha$ is reported with the zero-day results. Finally, the seed-level power floor is stated
rather than left implicit: with $n = 5$ seeds the smallest detectable paired effect is $d_z \approx 1.24$,
and $\approx 1.68$ at $80\%$ power, which is why an effect size accompanies every reported $p$-value.
.
% Requires: \usepackage{booktabs}, \usepackage{amssymb}, \usepackage{amsmath}.
% Cells marked \ddag are provisional (Day-14 dummy fidelity interface) and reprice on real prototypes.

\subsection{Data}
\label{sec:methods-data}

\paragraph{Datasets and unified representation.}
We evaluate on three public IoT/network intrusion-detection corpora: CIC-IoT2023, BoT-IoT and UNSW-NB15.
Two further corpora (TON\_IoT, Edge-IIoTset) were curated under the same pipeline but excluded, because
their held-out zero-day split does not survive the shared schema: projecting onto the common feature set
maps many distinct attack flows onto \emph{identical} vectors, and the cross-split de-duplication that
follows resolves collisions with training-split priority, so the surviving zero-day copies are dropped. The
held-out families end at one and two rows respectively, which cannot support a rejection-rate estimate; the
three retained datasets keep their zero-day splits intact. Since the three corpora publish incompatible
column vocabularies,
each is projected onto a shared \textbf{17-feature unified flow schema} (flow duration; total, source and
destination packet and byte counts; overall, source and destination rates; four packet-size statistics;
inter-arrival time; and harmonized \texttt{protocol} and \texttt{conn\_state} categoricals), together with a
fine multiclass label, a binary label, and a ten-family attack-ontology label.

\paragraph{Curation and leakage controls.}
Curation runs through a quantum-aware curation pipeline whose defining property is that it is
\emph{split-first}: every data-dependent parameter is fit on the training split alone and only applied to
the others. This departs deliberately from the conventional ordering, in which scaling or balancing precedes
splitting and thereby leaks evaluation data into fitted parameters. Concretely, the correlation prune
(dropping one of each $|r| \geq 0.95$ pair), the median/IQR robust scalers, the angle-encoding ranges
(mapping features to $[0, \pi]$ for angle embedding; the published benchmark uses the qubit budget of $8$
from the nested $4/8/12/16$ rankings), and the feature ranking are all fit on the training split and applied
unchanged to calibration, test and zero-day. Class balancing is applied to
training only; evaluation splits are never resampled.

Two further controls underpin the guarantee developed in Section~\ref{sec:methods-conformal}. First,
de-duplication is performed on the feature columns \emph{before} splitting, so one feature vector can land
in exactly one split; cross-split disjointness is then asserted on feature-content hashes rather than row
indices, which is what detects a repeated flow shared across splits. This is the relevant control for
CIC-IoT2023, whose leakage vector is repeated and near-duplicate flows rather than identifier columns.
Measured cross-split flow overlap is zero for every pair on every dataset, and within-split duplication is
zero. Second, the conformal threshold is calibrated on held-out \emph{known-class} data and never on the
unknown class: calibrating on the held-out family would leak precisely the signal the detector is meant to
discover.

\paragraph{Partition protocol.}
Each dataset is partitioned four ways --- Train / Calibration / Test / Zero-Day --- under a fixed seed
($42$), so the partitions are deterministic and reproduce exactly on re-run. The zero-day split is a
\emph{held-out attack family} rather than a random subsample: one or more entire attack classes are removed
from training, calibration and test and appear only at evaluation time. The remaining classes are stratified
$70/10/10/10$ into train/validation/calibration/test, with validation folded into train to give the
four-way conformal layout. Calibration contains known classes only, by construction. Three invariants are
asserted in the partitioner rather than merely documented: the zero-day family has zero rows outside the
zero-day split; calibration is a subset of the known classes; and the assignment is deterministic under the
seed.

\begin{table}[t]
\centering
\caption{Partition protocol. Row counts per split, number of known classes, and the held-out zero-day
family. Seed 42; deterministic on re-run.}
\label{tab:partitions}
\begin{tabular}{lrrrrrl}
\toprule
Dataset & Train & Calib. & Test & Zero-day & \#known & Held-out family \\
\midrule
CIC-IoT2023 & 151,049 & 18,883 & 18,883 & 10,984 & 31 & Mirai ($\times 3$) \\
BoT-IoT     & 148,729 & 18,591 & 18,591 &    683 &  4 & Theft \\
UNSW-NB15   &  81,215 & 10,112 & 10,086 &  1,216 &  8 & Shellcode, Worms \\
\bottomrule
\end{tabular}
\end{table}

\paragraph{Verifying exchangeability.}
Split-conformal validity requires \emph{exchangeability} between the calibration set and each test point ---
strictly weaker than i.i.d., and a condition that cannot be proven, only falsified. We therefore treat it as
an empirical claim and attempt to falsify it with a classifier two-sample test: a random forest is trained
to separate calibration rows from test rows and scored by five-fold cross-validated AUROC, where chance
performance means the two sets are indistinguishable on the features. We obtain AUROC $0.4999$, $0.5038$ and
$0.5035$ on CIC-IoT2023, BoT-IoT and UNSW-NB15 respectively, with cal-vs-test total-variation distances of
$0.0000$, $0.0000$ and $0.0019$; exchangeability is not rejected on any dataset. This claim is scoped to the
\emph{known-class} calibration--test split. The held-out zero-day split is deliberately non-exchangeable
with calibration --- that non-exchangeability is the detection signal --- so the guarantee below covers the
known-class portion of the test stream, and zero-day detection is reported separately.

\subsection{Conformal Statistics}
\label{sec:methods-conformal}

\paragraph{Nonconformity score and calibration.}
For a sample $x$ encoded as a density matrix $\rho_x$, the model supplies the Uhlmann fidelity
$F(\rho_x, \rho_c)$ to every known-class prototype $\rho_c$ (non-squared $F$). The nonconformity score is
\begin{equation}
s(x) \;=\; 1 - \max_{c} F(\rho_x, \rho_c),
\end{equation}
so a sample with low fidelity to every known prototype scores high and is more novel. Prototypes are built
from the training split only; were calibration rows to contribute to $\rho_c$, calibration scores would be
in-sample and the threshold optimistic. The threshold is the split-conformal order statistic of $s$ over the
known-class calibration set of size $n$,
\begin{equation}
k = \left\lceil (1-\alpha)(n+1) \right\rceil, \qquad q = s_{(k)},
\end{equation}
the $k$-th smallest calibration score. A test point is flagged zero-day iff $s > q$. Under exchangeability
of the known-class calibration and test points this yields the finite-sample marginal guarantee
\begin{equation}
1 - \alpha \;\leq\; \mathbb{P}(s_{\mathrm{test}} \leq q) \;\leq\; 1 - \alpha + \frac{1}{n+1},
\label{eq:coverage}
\end{equation}
i.e.\ the false-zero-day rate on known traffic is at most $\alpha$, with granularity $1/(n+1)$. Validity is
\emph{score-agnostic}: it does not require the fidelity score to be a good novelty signal. Detection
\emph{power} does, and is therefore reported as an empirical result rather than inherited from
Eq.~\eqref{eq:coverage} --- a coverage curve is not a power curve. We use one primary $\alpha = 0.05$ for
the quantum and classical arms alike, giving $k = 17{,}940$, $17{,}663$ and $9{,}608$ and thresholds
$q = 0.3006^{\ddag}$, $0.2629^{\ddag}$ and $0.3834^{\ddag}$ on the three datasets.

\paragraph{Marginal rather than class-conditional calibration.}
Class-conditional (Mondrian) thresholds are degenerate when a class holds fewer than
$\lceil 1/\alpha \rceil - 1$ calibration points, a floor of $19$ per class at $\alpha = 0.05$. CIC-IoT2023
has four known classes with at most nine calibration rows, so we calibrate marginally on all three datasets
for comparability and report class-conditional behaviour as a diagnostic. The scope of the marginal
guarantee is stated precisely, because it is easy to overclaim: marginal conformal controls the false-alarm
rate over the known-class \emph{mixture} and carries \emph{no per-class guarantee}. Rare classes contribute
few pooled calibration points and so do not skew $q$ --- the mixture guarantee is unaffected by class
imbalance --- but an individual class may still be over- or under-covered. Each known class's test flag
count is therefore judged against its own exact band; on the present interface, $0$ of $31$, $0$ of $4$ and
$0$ of $8$ classes fall outside their own band. Should a class fall outside on real prototypes, the
appropriate remedy is clustered conformal --- grouping rare classes until each cluster clears the
calibration floor --- rather than fully class-conditional conformal, which several classes cannot support.

\paragraph{Coverage verification.}
A bare $\mathrm{FZR} \leq \alpha$ assertion is the wrong acceptance test. Split conformal guarantees the
\emph{expectation}, marginally over the calibration and test draws; on a finite test set the empirical rate
fluctuates around $(n+1-k)/(n+1)$, so a naive threshold check rejects sound systems by luck roughly half the
time. We instead judge each observation against the exact finite-sample law: conditional on the calibration
draw, coverage is Beta-distributed, so the number of false flags $E$ on $m$ known test rows follows a
Beta-Binomial,
\begin{equation}
\mathrm{coverage} \mid \mathrm{cal} \sim \mathrm{Beta}(k,\, n{+}1{-}k)
\quad \Longrightarrow \quad
E \sim \mathrm{BetaBin}(m,\, n{+}1{-}k,\, k).
\end{equation}
For each (dataset, $\alpha$) we report the observed rate, its exact central $99\%$ band, and a one-sided
tail $p$-value $\mathbb{P}(E \geq e_{\mathrm{obs}})$. The band, not the point comparison against $\alpha$,
is the acceptance gate: a failure means the guarantee is genuinely violated rather than unlucky. Below-band
excursions are reported as conservativeness --- ties and score discreteness make the band conservative ---
not as failures. This is also the correct band for a calibration-draw plot; Clopper--Pearson describes a
fixed-threshold proportion and is the wrong object here. BoT-IoT illustrates the distinction: its achieved
rate $0.0530^{\ddag}$ exceeds $\alpha$ outright, yet sits comfortably inside the exact band
$[0.0443, 0.0559]$ at $p = 0.090$, which is ordinary sampling noise on $18{,}591$ test rows. Coverage is
verified across the full sweep $\alpha \in [0.01, 0.20]$ ($20$ values $\times$ $3$ datasets $=60$
verifications, all passing), and the achieved rate tracks the diagonal $y = \alpha$ inside the band across
the whole range.

\paragraph{Stability across seeds.}
Because the guarantee is marginal over the calibration/test draw, it must hold across draws rather than on
one split. We pool the known-class calibration and test scores, re-draw them at the original sizes under
seeds $42$--$46$, re-calibrate each re-split and judge it against its own exact band, holding the zero-day
pool fixed and scoring it at each re-split's $q$. All three datasets hold the exact band on all five
re-splits, and every $95\%$ confidence interval covers the target $\alpha = 0.05$ (mean achieved rates
$0.0517^{\ddag}$, $0.0501^{\ddag}$, $0.0508^{\ddag}$). The accompanying recall intervals quantify how little
the power number moves with the calibration draw.

\paragraph{Auditing the assumption at deployment.}
Coverage arithmetic is meaningful only if the assumption underneath it holds, so the deployed pipeline is
audited with six inferential tests per dataset --- two-sample Kolmogorov--Smirnov and Mann--Whitney on
calibration versus test scores, a class-mix $\chi^2$, a flag-rate-by-class $\chi^2$, and two index-drift
Spearman tests --- with Holm step-down over the declared family of $18$ cells. No violations are detected.
Two choices deserve statement. First, the textbook conformal $p$-value uniformity check is reported
descriptively and \emph{excluded} from the inferential family: every $p_i$ is computed against the same
calibration set, so the $p$-values are exchangeable rather than independent and the one-sample uniformity
null is far too tight (measured on this interface it rejects on approximately $37\%$ of genuinely
exchangeable re-splits); it is also redundant with the two-sample statistic up to $1/(n+1)$. Second, the
audit's power is demonstrated by injected-violation drills rather than assumed: trimming the calibration
tail, shifting test scores, and skewing the test class mix are each injected and each detected. The
class-mix injection is instructive --- the audit flags it decisively while the \emph{marginal} false-alarm
rate stays inside the band, which is exactly the failure a coverage-only check cannot see. A caveat travels
with the remedy: re-splitting repairs \emph{split-induced} non-exchangeability, and also makes a
\emph{deployment shift} look repaired, because calibration and test become exchangeable draws from the
shifted mixture --- but that mixture is not the deployment distribution. A genuine deployment shift is
answered by recalibrating on fresh data, never by re-splitting.

\paragraph{Significance protocol.}
All comparative claims pass through one pre-declared protocol, chosen so the test matches the object being
claimed. We report mean $\pm$ sample standard deviation (ddof $=1$) with $t$-based $95\%$ confidence
intervals. Seed-level comparisons use a two-sided paired $t$-test with pairs matched by seed over seeds
$42$--$46$; row-level comparisons use an exact two-sided binomial McNemar test on discordant pairs matched
by test row. Effect sizes accompany every $p$-value: paired $d_z = \mathrm{mean}(\Delta)/\mathrm{sd}(\Delta)$
for seed-level tests and Cohen's $h$ for proportions. Multiplicity is controlled by Holm--Bonferroni
step-down over a family declared \emph{before} testing, and a zero-variance guard forces $t = 0$, $p = 1$,
$d = 0$ so a self-comparison cannot fake significance. Seed-paired and row-paired tests answer different
questions and are reported separately rather than pooled: threshold-free ranking metrics speak to ordering,
whereas McNemar tests the flags a system actually deploys. A dry run on classical baselines showed the
correction earning its keep, with a raw $p = 0.011$ in a nine-test family correctly losing significance
under Holm at $p = 0.056$.

Families are declared per comparison surface: closed-set correctness against the classical multiclass
detector ($m = 3$, one comparison per dataset) and known-split false-alarm rate against every novelty head
($m = 8$; one dataset publishes no one-class SVM model, a real absence rather than a missing value),
corrected separately. We note explicitly that the false-alarm family compares systems at their
\emph{shipped operating points} --- the conformal detector at its guaranteed $\alpha = 0.05$, each
heuristic head at its own frozen threshold, which targeted a $2\alpha$ known-FPR budget --- so it supports a
deployed-behaviour statement rather than a matched-budget superiority claim; the like-for-like comparison at
one shared $\alpha$ is reported with the zero-day results. Finally, the seed-level power floor is stated
rather than left implicit: with $n = 5$ seeds the smallest detectable paired effect is $d_z \approx 1.24$,
and $\approx 1.68$ at $80\%$ power, which is why an effect size accompanies every reported $p$-value.
.
% Requires: \usepackage{booktabs}, \usepackage{amssymb}, \usepackage{amsmath}.
% Cells marked \ddag are provisional (Day-14 dummy fidelity interface) and reprice on real prototypes.

\subsection{Data}
\label{sec:methods-data}

\paragraph{Datasets and unified representation.}
We evaluate on three public IoT/network intrusion-detection corpora: CIC-IoT2023, BoT-IoT and UNSW-NB15.
Two further corpora (TON\_IoT, Edge-IIoTset) were curated under the same pipeline but excluded, because
their held-out zero-day split does not survive the shared schema: projecting onto the common feature set
maps many distinct attack flows onto \emph{identical} vectors, and the cross-split de-duplication that
follows resolves collisions with training-split priority, so the surviving zero-day copies are dropped. The
held-out families end at one and two rows respectively, which cannot support a rejection-rate estimate; the
three retained datasets keep their zero-day splits intact. Since the three corpora publish incompatible
column vocabularies,
each is projected onto a shared \textbf{17-feature unified flow schema} (flow duration; total, source and
destination packet and byte counts; overall, source and destination rates; four packet-size statistics;
inter-arrival time; and harmonized \texttt{protocol} and \texttt{conn\_state} categoricals), together with a
fine multiclass label, a binary label, and a ten-family attack-ontology label.

\paragraph{Curation and leakage controls.}
Curation runs through a quantum-aware curation pipeline whose defining property is that it is
\emph{split-first}: every data-dependent parameter is fit on the training split alone and only applied to
the others. This departs deliberately from the conventional ordering, in which scaling or balancing precedes
splitting and thereby leaks evaluation data into fitted parameters. Concretely, the correlation prune
(dropping one of each $|r| \geq 0.95$ pair), the median/IQR robust scalers, the angle-encoding ranges
(mapping features to $[0, \pi]$ for angle embedding; the published benchmark uses the qubit budget of $8$
from the nested $4/8/12/16$ rankings), and the feature ranking are all fit on the training split and applied
unchanged to calibration, test and zero-day. Class balancing is applied to
training only; evaluation splits are never resampled.

Two further controls underpin the guarantee developed in Section~\ref{sec:methods-conformal}. First,
de-duplication is performed on the feature columns \emph{before} splitting, so one feature vector can land
in exactly one split; cross-split disjointness is then asserted on feature-content hashes rather than row
indices, which is what detects a repeated flow shared across splits. This is the relevant control for
CIC-IoT2023, whose leakage vector is repeated and near-duplicate flows rather than identifier columns.
Measured cross-split flow overlap is zero for every pair on every dataset, and within-split duplication is
zero. Second, the conformal threshold is calibrated on held-out \emph{known-class} data and never on the
unknown class: calibrating on the held-out family would leak precisely the signal the detector is meant to
discover.

\paragraph{Partition protocol.}
Each dataset is partitioned four ways --- Train / Calibration / Test / Zero-Day --- under a fixed seed
($42$), so the partitions are deterministic and reproduce exactly on re-run. The zero-day split is a
\emph{held-out attack family} rather than a random subsample: one or more entire attack classes are removed
from training, calibration and test and appear only at evaluation time. The remaining classes are stratified
$70/10/10/10$ into train/validation/calibration/test, with validation folded into train to give the
four-way conformal layout. Calibration contains known classes only, by construction. Three invariants are
asserted in the partitioner rather than merely documented: the zero-day family has zero rows outside the
zero-day split; calibration is a subset of the known classes; and the assignment is deterministic under the
seed.

\begin{table}[t]
\centering
\caption{Partition protocol. Row counts per split, number of known classes, and the held-out zero-day
family. Seed 42; deterministic on re-run.}
\label{tab:partitions}
\begin{tabular}{lrrrrrl}
\toprule
Dataset & Train & Calib. & Test & Zero-day & \#known & Held-out family \\
\midrule
CIC-IoT2023 & 151,049 & 18,883 & 18,883 & 10,984 & 31 & Mirai ($\times 3$) \\
BoT-IoT     & 148,729 & 18,591 & 18,591 &    683 &  4 & Theft \\
UNSW-NB15   &  81,215 & 10,112 & 10,086 &  1,216 &  8 & Shellcode, Worms \\
\bottomrule
\end{tabular}
\end{table}

\paragraph{Verifying exchangeability.}
Split-conformal validity requires \emph{exchangeability} between the calibration set and each test point ---
strictly weaker than i.i.d., and a condition that cannot be proven, only falsified. We therefore treat it as
an empirical claim and attempt to falsify it with a classifier two-sample test: a random forest is trained
to separate calibration rows from test rows and scored by five-fold cross-validated AUROC, where chance
performance means the two sets are indistinguishable on the features. We obtain AUROC $0.4999$, $0.5038$ and
$0.5035$ on CIC-IoT2023, BoT-IoT and UNSW-NB15 respectively, with cal-vs-test total-variation distances of
$0.0000$, $0.0000$ and $0.0019$; exchangeability is not rejected on any dataset. This claim is scoped to the
\emph{known-class} calibration--test split. The held-out zero-day split is deliberately non-exchangeable
with calibration --- that non-exchangeability is the detection signal --- so the guarantee below covers the
known-class portion of the test stream, and zero-day detection is reported separately.

\subsection{Conformal Statistics}
\label{sec:methods-conformal}

\paragraph{Nonconformity score and calibration.}
For a sample $x$ encoded as a density matrix $\rho_x$, the model supplies the Uhlmann fidelity
$F(\rho_x, \rho_c)$ to every known-class prototype $\rho_c$ (non-squared $F$). The nonconformity score is
\begin{equation}
s(x) \;=\; 1 - \max_{c} F(\rho_x, \rho_c),
\end{equation}
so a sample with low fidelity to every known prototype scores high and is more novel. Prototypes are built
from the training split only; were calibration rows to contribute to $\rho_c$, calibration scores would be
in-sample and the threshold optimistic. The threshold is the split-conformal order statistic of $s$ over the
known-class calibration set of size $n$,
\begin{equation}
k = \left\lceil (1-\alpha)(n+1) \right\rceil, \qquad q = s_{(k)},
\end{equation}
the $k$-th smallest calibration score. A test point is flagged zero-day iff $s > q$. Under exchangeability
of the known-class calibration and test points this yields the finite-sample marginal guarantee
\begin{equation}
1 - \alpha \;\leq\; \mathbb{P}(s_{\mathrm{test}} \leq q) \;\leq\; 1 - \alpha + \frac{1}{n+1},
\label{eq:coverage}
\end{equation}
i.e.\ the false-zero-day rate on known traffic is at most $\alpha$, with granularity $1/(n+1)$. Validity is
\emph{score-agnostic}: it does not require the fidelity score to be a good novelty signal. Detection
\emph{power} does, and is therefore reported as an empirical result rather than inherited from
Eq.~\eqref{eq:coverage} --- a coverage curve is not a power curve. We use one primary $\alpha = 0.05$ for
the quantum and classical arms alike, giving $k = 17{,}940$, $17{,}663$ and $9{,}608$ and thresholds
$q = 0.3006^{\ddag}$, $0.2629^{\ddag}$ and $0.3834^{\ddag}$ on the three datasets.

\paragraph{Marginal rather than class-conditional calibration.}
Class-conditional (Mondrian) thresholds are degenerate when a class holds fewer than
$\lceil 1/\alpha \rceil - 1$ calibration points, a floor of $19$ per class at $\alpha = 0.05$. CIC-IoT2023
has four known classes with at most nine calibration rows, so we calibrate marginally on all three datasets
for comparability and report class-conditional behaviour as a diagnostic. The scope of the marginal
guarantee is stated precisely, because it is easy to overclaim: marginal conformal controls the false-alarm
rate over the known-class \emph{mixture} and carries \emph{no per-class guarantee}. Rare classes contribute
few pooled calibration points and so do not skew $q$ --- the mixture guarantee is unaffected by class
imbalance --- but an individual class may still be over- or under-covered. Each known class's test flag
count is therefore judged against its own exact band; on the present interface, $0$ of $31$, $0$ of $4$ and
$0$ of $8$ classes fall outside their own band. Should a class fall outside on real prototypes, the
appropriate remedy is clustered conformal --- grouping rare classes until each cluster clears the
calibration floor --- rather than fully class-conditional conformal, which several classes cannot support.

\paragraph{Coverage verification.}
A bare $\mathrm{FZR} \leq \alpha$ assertion is the wrong acceptance test. Split conformal guarantees the
\emph{expectation}, marginally over the calibration and test draws; on a finite test set the empirical rate
fluctuates around $(n+1-k)/(n+1)$, so a naive threshold check rejects sound systems by luck roughly half the
time. We instead judge each observation against the exact finite-sample law: conditional on the calibration
draw, coverage is Beta-distributed, so the number of false flags $E$ on $m$ known test rows follows a
Beta-Binomial,
\begin{equation}
\mathrm{coverage} \mid \mathrm{cal} \sim \mathrm{Beta}(k,\, n{+}1{-}k)
\quad \Longrightarrow \quad
E \sim \mathrm{BetaBin}(m,\, n{+}1{-}k,\, k).
\end{equation}
For each (dataset, $\alpha$) we report the observed rate, its exact central $99\%$ band, and a one-sided
tail $p$-value $\mathbb{P}(E \geq e_{\mathrm{obs}})$. The band, not the point comparison against $\alpha$,
is the acceptance gate: a failure means the guarantee is genuinely violated rather than unlucky. Below-band
excursions are reported as conservativeness --- ties and score discreteness make the band conservative ---
not as failures. This is also the correct band for a calibration-draw plot; Clopper--Pearson describes a
fixed-threshold proportion and is the wrong object here. BoT-IoT illustrates the distinction: its achieved
rate $0.0530^{\ddag}$ exceeds $\alpha$ outright, yet sits comfortably inside the exact band
$[0.0443, 0.0559]$ at $p = 0.090$, which is ordinary sampling noise on $18{,}591$ test rows. Coverage is
verified across the full sweep $\alpha \in [0.01, 0.20]$ ($20$ values $\times$ $3$ datasets $=60$
verifications, all passing), and the achieved rate tracks the diagonal $y = \alpha$ inside the band across
the whole range.

\paragraph{Stability across seeds.}
Because the guarantee is marginal over the calibration/test draw, it must hold across draws rather than on
one split. We pool the known-class calibration and test scores, re-draw them at the original sizes under
seeds $42$--$46$, re-calibrate each re-split and judge it against its own exact band, holding the zero-day
pool fixed and scoring it at each re-split's $q$. All three datasets hold the exact band on all five
re-splits, and every $95\%$ confidence interval covers the target $\alpha = 0.05$ (mean achieved rates
$0.0517^{\ddag}$, $0.0501^{\ddag}$, $0.0508^{\ddag}$). The accompanying recall intervals quantify how little
the power number moves with the calibration draw.

\paragraph{Auditing the assumption at deployment.}
Coverage arithmetic is meaningful only if the assumption underneath it holds, so the deployed pipeline is
audited with six inferential tests per dataset --- two-sample Kolmogorov--Smirnov and Mann--Whitney on
calibration versus test scores, a class-mix $\chi^2$, a flag-rate-by-class $\chi^2$, and two index-drift
Spearman tests --- with Holm step-down over the declared family of $18$ cells. No violations are detected.
Two choices deserve statement. First, the textbook conformal $p$-value uniformity check is reported
descriptively and \emph{excluded} from the inferential family: every $p_i$ is computed against the same
calibration set, so the $p$-values are exchangeable rather than independent and the one-sample uniformity
null is far too tight (measured on this interface it rejects on approximately $37\%$ of genuinely
exchangeable re-splits); it is also redundant with the two-sample statistic up to $1/(n+1)$. Second, the
audit's power is demonstrated by injected-violation drills rather than assumed: trimming the calibration
tail, shifting test scores, and skewing the test class mix are each injected and each detected. The
class-mix injection is instructive --- the audit flags it decisively while the \emph{marginal} false-alarm
rate stays inside the band, which is exactly the failure a coverage-only check cannot see. A caveat travels
with the remedy: re-splitting repairs \emph{split-induced} non-exchangeability, and also makes a
\emph{deployment shift} look repaired, because calibration and test become exchangeable draws from the
shifted mixture --- but that mixture is not the deployment distribution. A genuine deployment shift is
answered by recalibrating on fresh data, never by re-splitting.

\paragraph{Significance protocol.}
All comparative claims pass through one pre-declared protocol, chosen so the test matches the object being
claimed. We report mean $\pm$ sample standard deviation (ddof $=1$) with $t$-based $95\%$ confidence
intervals. Seed-level comparisons use a two-sided paired $t$-test with pairs matched by seed over seeds
$42$--$46$; row-level comparisons use an exact two-sided binomial McNemar test on discordant pairs matched
by test row. Effect sizes accompany every $p$-value: paired $d_z = \mathrm{mean}(\Delta)/\mathrm{sd}(\Delta)$
for seed-level tests and Cohen's $h$ for proportions. Multiplicity is controlled by Holm--Bonferroni
step-down over a family declared \emph{before} testing, and a zero-variance guard forces $t = 0$, $p = 1$,
$d = 0$ so a self-comparison cannot fake significance. Seed-paired and row-paired tests answer different
questions and are reported separately rather than pooled: threshold-free ranking metrics speak to ordering,
whereas McNemar tests the flags a system actually deploys. A dry run on classical baselines showed the
correction earning its keep, with a raw $p = 0.011$ in a nine-test family correctly losing significance
under Holm at $p = 0.056$.

Families are declared per comparison surface: closed-set correctness against the classical multiclass
detector ($m = 3$, one comparison per dataset) and known-split false-alarm rate against every novelty head
($m = 8$; one dataset publishes no one-class SVM model, a real absence rather than a missing value),
corrected separately. We note explicitly that the false-alarm family compares systems at their
\emph{shipped operating points} --- the conformal detector at its guaranteed $\alpha = 0.05$, each
heuristic head at its own frozen threshold, which targeted a $2\alpha$ known-FPR budget --- so it supports a
deployed-behaviour statement rather than a matched-budget superiority claim; the like-for-like comparison at
one shared $\alpha$ is reported with the zero-day results. Finally, the seed-level power floor is stated
rather than left implicit: with $n = 5$ seeds the smallest detectable paired effect is $d_z \approx 1.24$,
and $\approx 1.68$ at $80\%$ power, which is why an effect size accompanies every reported $p$-value.
